\section*{Capítulo \ref{ch:b2:biogeochemical-cycles} --- Ciclos biogeoquímicos}

\begin{solution}{exo:b2:biogeochemical-cycles:1}
Un \omterm{def:b2:biogeochemical-cycles:reservoir}{reservorio} es una reserva del elemento (una masa); un \omterm{def:b2:biogeochemical-cycles:reservoir}{flujo} es una
velocidad de transferencia entre \omterm{def:b2:biogeochemical-cycles:reservoir}{reservorios} (masa por año); el \omterm{def:b2:biogeochemical-cycles:reservoir}{tiempo de
residencia} es la masa del \omterm{def:b2:biogeochemical-cycles:reservoir}{reservorio} dividida por su salida total, el
tiempo medio que pasa un átomo en él. Plantas terrestres: $450/120 \approx 4$
años (el carbono de una hoja, aproximadamente; la \omterm{def:b2:plant-meristems:wood}{madera} lo retiene
durante décadas y la media oculta la dispersión). Océano profundo:
$37000/100 \approx 400$ años --- el tiempo que tarda una molécula que se ha
hundido en volver a ver la superficie.
\end{solution}

\begin{solution}{exo:b2:biogeochemical-cycles:2}
$2.5\times 2.13 = \qty{5.3}{GtC}$ al año. $10/2.13 = \qty{4.7}{ppm}$
--- si todo se quedara en el aire; con una \omterm{prop:b2:biogeochemical-cycles:carbon}{fracción atmosférica} del
\qty{45}{\%}, \qty{2.1}{ppm}.
\end{solution}

\begin{solution}{exo:b2:biogeochemical-cycles:3}
Apertura: la fijación biológica por la nitrogenasa (cianobacterias,
rizobios, \emph{Azotobacter}, \emph{Frankia}) y los rayos. Retorno: la
\omterm{def:b2:microbial-metabolism:anaerobic}{desnitrificación} (anaerobios facultativos como \emph{Pseudomonas}
y \emph{Paracoccus} en suelos y sedimentos anóxicos) y el anammox
(planctomicetos). Entre ambos: la amonificación del nitrógeno orgánico a
amonio por los descomponedores (hongos, bacterias); la nitrificación del
amonio a nitrito (\emph{Nitrosomonas}, arqueas oxidantes del amoníaco) y
del nitrito a nitrato (\emph{Nitrobacter}); y la asimilación del amonio o
del nitrato por las plantas y los microorganismos.
\end{solution}

\begin{solution}{exo:b2:biogeochemical-cycles:4}
A escala geológica, el nitrógeno siempre puede obtenerse del inagotable
$\mathrm{N_2}$ del aire gracias a los fijadores, que resultan favorecidos
siempre que falta nitrógeno; el fósforo no tiene un \omterm{def:b2:biogeochemical-cycles:reservoir}{reservorio} así y su
aporte lo fija la meteorización, de modo que el techo de la productividad
a largo plazo es el fósforo. Dentro de una estación, la fijación es lenta
y costosa, y la mayoría de las plantas no pueden realizarla, de modo que
el nitrógeno disponible es lo que limita el crecimiento aquí y ahora ---
y por eso los abonos son sobre todo nitrogenados.
\end{solution}

\begin{solution}{exo:b2:biogeochemical-cycles:5}
$\tau = 1/0.2 = 5$ años; $M^{*} = F\tau = \qty{10}{t}$, es decir,
$\qty{1}{g/m^3} = \qty{1000}{mg/m^3}$, extremadamente eutrófico. Tiempo
hasta el \qty{90}{\%}: $\tau\ln 10 = 11.5$ años. Reducir el aporte a la
mitad reduce a la mitad el \omterm{def:b2:biogeochemical-cycles:reservoir}{estado estacionario}, hasta \qty{500}{mg/m^3},
que se alcanza en la misma escala de tiempo.
\end{solution}

\begin{solution}{exo:b2:biogeochemical-cycles:6}
Década de 1960: \qty{0.9}{ppm/yr}, \qty{1.9}{GtC/yr}. Década de 2010:
\qty{2.3}{ppm/yr}, \qty{5.0}{GtC/yr}. La tasa de crecimiento se ha
multiplicado por dos y medio, al mismo ritmo que las emisiones.
\end{solution}

\begin{solution}{exo:b2:biogeochemical-cycles:7}
\qty{6}{ppm} son \qty{13}{GtC}, frente a una producción primaria neta de
\qty{35}{GtC}: los dientes de sierra equivalen a un tercio. Son menores
porque la respiración y la descomposición continúan durante el verano (los
dientes de sierra registran la absorción neta, no la bruta), porque las
estaciones del hemisferio sur son opuestas y compensan en parte, y porque
el océano amortigua la oscilación.
\end{solution}

\begin{solution}{exo:b2:biogeochemical-cycles:8}
$\qty{200}{kg}/(\qty{28}{g/mol}) = 7100$~mol de $\mathrm{N_2}$; $16\times
7100 = \num{114000}$~mol de ATP; $/30 = 3800$~mol de glucosa, \qty{690}{kg}.
Un cultivo que produce \qty{10}{t} de materia seca fija, contando la
respiración, unas
\qtyrange{15}{20}{t} de azúcar, de modo que la fijación le cuesta alrededor
del \qty{4}{\%} de sus fotosintatos --- el precio de independizarse del
nitrógeno del suelo.
\end{solution}

\begin{solution}{exo:b2:biogeochemical-cycles:9}
Cociente disponible $\mathrm{N:P} = 32/2.2 = 14.5$, por debajo del 16 de
Redfield: el nitrato se agota primero, y sobran \qty{0.2}{\micro mol/kg}
de fosfato. Carbono fijado: $32\times 106/16 =
\qty{212}{\micro mol/kg}$, unos \qty{2.5}{mg} de carbono por
kilogramo de agua.
\end{solution}

\begin{solution}{exo:b2:biogeochemical-cycles:10}
$E^{*} = Q\tau = \qty{500}{GtC}$, \qty{235}{ppm} por encima de los 285
preindustriales: unas \qty{520}{ppm}. La realidad es peor porque los
sumideros no son un único compartimento rápido: el océano superficial se
satura (su capacidad de absorber $\mathrm{CO_2}$ disminuye a medida que se
consume su carbonato), el sumidero terrestre se debilita con el
calentamiento, y la eliminación real tiene un componente lento --- la mezcla
con el océano profundo y la meteorización de las rocas --- con constantes
de tiempo de mil a cien mil años, de modo que una fracción del exceso no se
relaja con ningún $\tau$ de décadas.
\end{solution}

\begin{solution}{exo:b2:biogeochemical-cycles:11}
Lo que se midió es el intercambio del $\mathrm{CO_2}$ atmosférico con el
océano y la biosfera, no la desintegración: el $\mathrm{{}^{14}C}$ de las bombas se
diluyó en los compartimentos, mucho mayores, del océano superficial y de
las plantas y los suelos, que devolvieron carbono ordinario. La semivida
de 11 años del exceso (un tiempo de e-plegado de unos 16 años) es la escala
de tiempo de ese intercambio --- la renovación del carbono de la atmósfera
con respecto al océano superficial y a las tierras emergidas, más larga que
el \omterm{def:b2:biogeochemical-cycles:reservoir}{tiempo de residencia} bruto de cuatro años porque buena parte de lo que
sale vuelve enseguida.
\end{solution}

\begin{solution}{exo:b2:biogeochemical-cycles:12}
Nitrógeno: fijación natural de unos 200~Tg al año; fijación humana
(Haber--Bosch 120, cultivos 60, combustión 30) de unos 210: el \omterm{def:b2:biogeochemical-cycles:reservoir}{flujo} se ha
duplicado, y los \omterm{def:b2:biogeochemical-cycles:reservoir}{reservorios} afectados --- suelos, ríos, mares costeros, el
$\mathrm{N_2O}$ del aire --- son pequeños y rápidos, de modo que el cambio
se nota en décadas, mientras que el \omterm{def:b2:biogeochemical-cycles:reservoir}{reservorio} de $\mathrm{N_2}$ no se toca.
Carbono: las emisiones humanas de 11~GtC al año son solo el \qty{5}{\%} del
intercambio natural bruto, de unas 200, de modo que el \omterm{def:b2:biogeochemical-cycles:reservoir}{flujo} apenas se
roza; pero los \omterm{def:b2:biogeochemical-cycles:reservoir}{flujos} naturales se anulan y el humano no, y ha elevado el
\omterm{def:b2:biogeochemical-cycles:reservoir}{reservorio} atmosférico un \qty{50}{\%} --- el pequeño roce tiene un gran
efecto acumulado sobre un \omterm{def:b2:biogeochemical-cycles:reservoir}{reservorio} con un sumidero último lento. Qué
ciclo está ``más perturbado'' depende de si se cuentan los \omterm{def:b2:biogeochemical-cycles:reservoir}{flujos} o la
acumulación.
\end{solution}

\begin{solution}{pb:b2:biogeochemical-cycles:1}
\textbf{1.} $285\times 2.13 = \qty{607}{GtC}$; $412\times 2.13 =
\qty{878}{GtC}$; aumento de \qty{271}{GtC}.
\textbf{2.} $271/700 = 0.39$ (la fracción de las últimas décadas, contando
las emisiones por cambio de uso del suelo, se acerca más a $0.45$).
\textbf{3.} El océano, alrededor de una cuarta parte de las emisiones, y
las tierras emergidas (rebrote de los bosques y fertilización por
$\mathrm{CO_2}$), alrededor de un tercio; el resto de las
\qty{430}{GtC} que no están en el aire.
\textbf{4.} $175/38000 = \qty{0.5}{\%}$ --- despreciable para el océano
entero; pero la absorción ha entrado hasta ahora sobre todo en la capa
superficial (\qty{900}{GtC}), donde supone un aumento de una décima parte o
más, y la química de los carbonatos convierte un pequeño aumento del
$\mathrm{CO_2}$ disuelto en un aumento mayor de la concentración de iones
hidrógeno: el pH superficial ha bajado
$0.1$, un aumento de la acidez del \qty{30}{\%}.
\textbf{5.} Con $Q = Q_0 e^{t/T}$, se prueba $E = Ae^{t/T}$: $A/T = Q_0 -
A/\tau$, de modo que $A = Q_0\tau T/(\tau + T)$ y la \omterm{prop:b2:biogeochemical-cycles:carbon}{fracción atmosférica}
$\dot E/Q = \tau/(\tau + T)$ es constante. Unas emisiones que crecen un
\qty{2}{\%} al año ($T = 50$) con $\tau = 40$ dan $0.44$: la fracción es
constante porque el sumidero y la fuente crecen juntos.
\textbf{6.} Con $Q$ constante, $E \to Q\tau$ y $\dot E \to 0$: la
\omterm{prop:b2:biogeochemical-cycles:carbon}{fracción atmosférica} tiende a cero a medida que los sumideros se ponen al
día --- en el modelo de una caja; en realidad los sumideros se saturan y la
fracción baja más despacio.
\textbf{7.} \qty{271}{GtC}.
\textbf{8.} $E^{*} = 10\times 100 = \qty{1000}{GtC}$: $285 +
1000/2.13 = \qty{755}{ppm}$.
\textbf{9.} $E(t) = 1000 + (271 - 1000)e^{-t/100}$; con $t = 80$,
$E = 1000 - 729\times 0.449 = \qty{672}{GtC}$: \qty{600}{ppm}.
\textbf{10.} $Q = 10(1 - t/50)$: solución particular $E_p = at + b$
con $a = -20$, $b = 3000$ (de $a = 10 - b/\tau$ y $0 = -0.2 -
a/\tau$); $E = 3000 - 20t + (271 - 3000)e^{-t/100}$. Con $t = 50$:
$2000 - 2729\times 0.607 = \qty{344}{GtC}$, \qty{447}{ppm}. Después, $E$
decae libremente: con $t = 80$, $344\,e^{-0.3} = \qty{255}{GtC}$,
\qty{405}{ppm}.
\textbf{11.} \qty{600}{ppm} frente a \qty{405}{ppm}: el segundo escenario
devuelve el aire, en 2100, casi a donde está hoy.
\textbf{12.} El $\tau$ único hace que los sumideros sigan trabajando al
mismo ritmo sobre un exceso que disminuye; en realidad, los sumideros
rápidos (la capa de mezcla, los bosques en crecimiento) ya han absorbido
lo que pueden, y la eliminación restante se debe a la lenta mezcla con el
océano profundo y a la meteorización, con $\tau$ de siglos a milenios, de
modo que el descenso tras alcanzar las emisiones netas nulas es mucho más
lento --- la concentración se mantiene aproximadamente constante durante
siglos en lugar de bajar.
\textbf{13.} Cultivo, \qty{90}{kg}; desnitrificado, \qty{37.5}{kg};
lixiviado, \qty{22.5}{kg} de nitrógeno por hectárea y por año.
\textbf{14.} $\qty{22.5}{kg}/\qty{14}{g/mol} = 1600$~mol de N; carbono:
$1600\times 106/16 = \num{10600}$~mol, \qty{128}{kg}.
\textbf{15.} \num{10600}~mol de $\mathrm{O_2}$, \qty{340}{kg}.
\textbf{16.} $\num{10600}/0.25 = \num{42000}$~\unit{m^3}: la lixiviación
de una hectárea puede quitar el oxígeno a cuarenta mil metros cúbicos de
agua de fondo cada año --- y el Misisipi drena cien millones de hectáreas.
\textbf{17.} $37.5\times 0.01 = \qty{0.375}{kg}$ de nitrógeno en forma de
$\mathrm{N_2O}$, es decir, $0.375\times 44/28 = \qty{0.59}{kg}$ de $\mathrm{N_2O}$;
$\times 270 = \qty{160}{kg}$ de $\mathrm{CO_2}$ equivalente.
\textbf{18.} $150\times 4 = \qty{600}{kg}$ de $\mathrm{CO_2}$: la fabricación
pesa cuatro veces más que el $\mathrm{N_2O}$ del campo, y ambos juntos suman
tres cuartos de tonelada por hectárea y por año.
\textbf{19.} $\tau = 10$ años; $M^{*} = 5\times 10 = \qty{50}{t}$;
concentración: $50/(5\times 10^{7}) = \qty{1}{g/m^3} =
\qty{1000}{mg/m^3}$.
\textbf{20.} Treinta veces el umbral: gravemente eutrófico.
\textbf{21.} $1/\tau = 0.1 + 0.2 = 0.3$: $\tau = 3.3$ años; $M^{*} =
5/0.3 = \qty{16.7}{t}$, \qty{330}{mg/m^3}.
\textbf{22.} $M^{*} = 1/0.3 = \qty{3.3}{t}$, \qty{67}{mg/m^3} --- todavía
eutrófico; a menos del \qty{10}{\%} al cabo de $\tau\ln 10 = 7.7$ años.
\textbf{23.} El sedimento no es un sumidero de un solo sentido: el fósforo
almacenado en él durante los años eutróficos se libera de nuevo en las
condiciones anóxicas que crea la propia eutrofización (el fosfato unido al
hierro se disuelve cuando el hierro se reduce), de modo que el lago se
alimenta a sí mismo durante años después de cortar el aporte externo ---
la carga interna, un segundo \omterm{def:b2:biogeochemical-cycles:reservoir}{reservorio} con un \omterm{def:b2:biogeochemical-cycles:reservoir}{tiempo de residencia} largo.
\textbf{24.} $\qty{5000}{kg}/\qty{31}{g/mol} = 1.6\times 10^{5}$~mol de P;
carbono: $\times 106 = 1.7\times 10^{7}$~mol, \qty{205}{t} de carbono al
año --- unos \qty{20}{g/m^2} sobre la superficie del lago si tiene
\qty{5}{m} de profundidad, una proliferación intensa.
\textbf{25.} \omterm{prop:b2:biogeochemical-cycles:carbon}{Fracción atmosférica} de $0.39$ durante 1850--2020; en 2100,
\qty{600}{ppm} con las emisiones mantenidas en \qty{10}{GtC/yr} y
\qty{405}{ppm} con las emisiones reducidas a cero en 2070 (un modelo
optimista de $\tau$ único); nitrógeno lixiviado, \qty{22.5}{kg} por
hectárea y por año.
\end{solution}
